\Needspace{5\baselineskip}\textbf{Problems}

\begin{enumerate}[label=\thesection.\arabic*.,ref=\thesection.\arabic*,start=17,series=book-problems]

\item\label{ex:orthogonal-projection-spectrum} 描述 Example~\ref{ex:orthogonal-projection-operator} 中所述投影算子的谱 $\sigma(P_M)$.

\item\label{ex:continuous-multiplication-operator-spectrum} \begin{enumerate}[label=(\alph*)]\item 对每个 $u\in C_b([0,1])$, 用下式定义乘法算子 $M:C_b([0,1])\to C_b([0,1])$:\[M(u)(x):=xu(x),\]描述 $\sigma(M)$.\item 给定 $v\in C_b([0,1])$. 对每个 $u\in C_b([0,1])$, 用下式定义算子 $M_v:C_b([0,1])\to C_b([0,1])$:\[M_v(u)(x):=v(x)u(x),\]描述 $\sigma(M_v)$.\end{enumerate}

\item\label{ex:bounded-operator-extension-spectral-inclusions} 假设 $(\mathcal{D}(\widetilde A),\widetilde A)$ 是有界算子 $(\mathcal{D}(A),A)$ 的扩张. 证明以下结果:\begin{enumerate}[label=(\alph*)]
\item $\sigma_p(\widetilde A)\supset\sigma_p(A)$.
\item $\sigma_r(\widetilde A)\subset\sigma_r(A)$.
\item $\sigma_c(A)\subset\sigma_c(\widetilde A)\cup\sigma_p(\widetilde A)$.
\item $\rho(\widetilde A)\subset\rho(A)\cup\sigma_r(A)$.
\end{enumerate}

\item\label{ex:bounded-operator-resolvent-decay-at-infinity} 设 $A\in\mathcal{L}(X)$. 证明当 $|\lambda|\to\infty$ 时, $\|R_\lambda(A)\|\to0$.

\item\label{ex:unbounded-second-derivative-noncompact-spectrum} 设\[\mathcal{D}(A)=\{u\in H^2(0,1)\mid u(0)=u(1)=0\}.\]对于 $u\in\mathcal{D}(A)$, 用下式定义从 $L^2(0,1)$ 到 $L^2(0,1)$ 的算子 $(\mathcal{D}(A),A)$:\[Au=u''\]证明 $\sigma(A)$ 不是紧集. 你的答案是否与 Corollary~\ref{cor:bounded-operator-spectrum-compact} 矛盾?

\item\label{ex:weak-and-strong-banach-valued-analyticity} 设 $G\subset\mathbb{C}$ 为区域, $X$ 为 Banach 空间. 对于映射
\begin{equation}\label{eq:weakly-analytic-banach-valued-map}
\mathbb{C}\supset G\ni\lambda\mapsto B(\lambda)\in X
\end{equation}
若对每个 $g\in X^*$, 由下式定义的复值函数
\begin{equation}\label{eq:weakly-analytic-scalar-functional-composition}
f(\lambda):=g(B(\lambda))
\end{equation}
都在 $\lambda_0$ 的某个邻域内 (按通常意义) 解析, 则称它在 $\lambda_0\in\mathbb{C}$ 处\textbf{弱解析}. 若函数 $B(\lambda)$ 在 $G$ 中每一点都解析, 则称它在 $G$ 上解析.\begin{enumerate}[label=(\alph*)]\item 证明 (强) 解析性蕴含弱解析性.\item 证明弱解析性蕴含 (强) 解析性.\end{enumerate}

\item\label{ex:closed-unbounded-resolvent-analyticity} 将 Theorems~\ref{thm:local-resolvent-power-series} 和 \ref{thm:bounded-operator-analytic-resolvent} 推广到无界闭算子.

\end{enumerate}
